% ------------------------------------------------------------
% Chapter 8: Cosmic Microwave Background
% ------------------------------------------------------------

\chapter{Cosmic Microwave Background}
\label{chap:cmb}

\section{Overview}

The cosmic microwave background (CMB) provides one of the most precise tests of
cosmological models. In the \rsvpabbr{} framework, temperature anisotropies arise
from both gravitational potential fluctuations and lamphrodynamic entropy
inhomogeneities. The goal of this chapter is to derive the angular power spectrum
$C_\ell$ in \rsvpabbr{}, and to compare with observational data, notably the
\textsc{Planck} 2018 results.

\section{Boltzmann Equation for CMB Photons}

CMB photons propagate through the lamphrodynamic plenum according to a
modified Boltzmann equation,
\begin{equation}
  \label{eq:boltzmann-rsvp}
  \frac{df}{d\lambda}
   = \mathcal{C}[f] + \mathcal{L}_{\mathrm{lam}}[f],
\end{equation}
where $\mathcal{C}[f]$ denotes collision terms (e.g.\ Thomson scattering) and
$\mathcal{L}_{\mathrm{lam}}[f]$ encodes lamphrodynamic effects due to
$(\PhiField,\vField,\SField)$. For scalar perturbations,
$\mathcal{L}_{\mathrm{lam}}$ contributes additional potential and entropy
terms influencing photon trajectories.

\section{Sachs--Wolfe and Integrated Sachs--Wolfe Effects}

Anisotropies in the CMB temperature receive contributions from the
Sachs--Wolfe (SW) and integrated Sachs--Wolfe (ISW) effects. In
\rsvpabbr{},
\begin{equation}
  \label{eq:sachs-wolfe-rsvp}
  \frac{\delta T}{T}\Big|_{\mathrm{SW}}
  = \frac{1}{3}\Phi_{\mathrm{eff}},
\end{equation}
where $\Phi_{\mathrm{eff}}$ includes lamphrodynamic corrections to the
gravitational potential. The ISW effect arises from time variation of
$\Phi_{\mathrm{eff}}$ along the line of sight. A detailed derivation yields
\begin{equation}
  \label{eq:isw-rsvp}
  \left(\frac{\delta T}{T}\right)_{\mathrm{ISW}}
  = \int_{\text{LOS}}\dot{\Phi}_{\mathrm{eff}}\,d\eta,
\end{equation}
which depends on the lamphrodynamic source terms $\SField$ and $\vField$.

\section{Angular Power Spectrum}

The angular power spectrum in \rsvpabbr{} is defined by
\[
  C_\ell
  := \frac{1}{2\ell+1}\sum_{m=-\ell}^\ell
       |\langle a_{\ell m}\rangle|^2,
\]
where $a_{\ell m}$ are the spherical–harmonic coefficients of the 
temperature anisotropy. The lamphrodynamic enhancements alter the 
transfer functions, leading to modifications of the acoustic peak
structure and Silk damping tail.

\begin{remark}[Peak shifts]
Lamphrodynamic contributions can shift both peak positions and amplitudes,
depending on the scale--dependence of $\rho_{\mathrm{lam}}$ and
$p_{\mathrm{lam}}$ during recombination.
\end{remark}

\section{Comparison with Planck Data}

Matching the \rsvpabbr{} $C_\ell$ with \textsc{Planck} 2018 data provides
stringent constraints on lamphrodynamic parameters. A least--squares fit in
the space of $(\alpha_i,\beta_i)$ yields best--fit constraints and
confidence intervals. Preliminary analysis indicates that certain classes of
lamphrodynamic models can match current data, but full analysis requires
numerical Boltzmann solvers (see \cref{app:cosmo-numerics}).

\section{Discussion}

The CMB provides some of the strongest constraints on \rsvpabbr{}.
Matching the observed acoustic peaks and damping tail is challenging but
feasible for suitable lamphrodynamic couplings. Further development of
numerical Boltzmann codes tailored to \rsvpabbr{} is needed for definitive
conclusions.

\clearpage
